\section{真实照片结果与几何质量}
\subsection{预处理比较}
表~\ref{tab:real-filters} 固定三层与 NMS，汇总十一图重复率。高斯低通并非所有方法的最优设置：LoG 在该设置下为 $0.749$，Sobel 的高斯高通为 $0.776$，而 SIFT 不滤波为 $0.631$、理想低通为 $0.634$，均高于其高斯低通的 $0.547$。这与印刷字细节、响应类型和描述子不同有关，但单一书页不能支持普遍排名。

\begin{table}[H]
\centering\small
\caption{预处理与管线存在交互。三层、NMS 开启时十一图平均重复率，来自主矩阵三层子集；不是按匹配总数池化的比例。}
\label{tab:real-filters}
\begin{tabular}{lrrrrr}
\toprule
方法 & 不滤波 & 理想低通 & 高斯低通 & 理想高通 & 高斯高通\\
\midrule
Sobel & 0.739 & 0.647 & 0.716 & 0.725 & 0.776\\
LoG & 0.738 & 0.535 & 0.749 & 0.662 & 0.672\\
Harris & 0.686 & 0.520 & 0.709 & 0.684 & 0.664\\
Shi-Tomasi & 0.653 & 0.525 & 0.706 & 0.651 & 0.651\\
SIFT & 0.631 & 0.634 & 0.547 & 0.403 & 0.367\\
\bottomrule
\end{tabular}
\end{table}

\subsection{完整设置：比例不能替代数量}
表~\ref{tab:real-full} 是高斯低通、三层与 NMS 的共同完整设置，不是逐方法挑选最好配置。Sobel 重复率 $0.716$，但每图平均接受匹配只有 $1.91$，正确匹配 $0.73$。SIFT 平均接受 $4.91$、正确 $2.73$，也不能仅凭这些点数保证稳定的几何估计；本实验的 $H$ 是由标注另行拟合，而不是通过这些稀少匹配估计。

此外，SIFT 在不滤波消融中正确率均值为 $0.923$，平均正确匹配 $5.45$，高于完整设置的 $0.541$ 和 $2.73$。该现象说明，对这些处理后的 JPEG，额外高斯低通可能损失 SIFT 所需的细节，不应把“开启滤波”默认等同于更可信。

\begin{table}[H]
\centering\small
\caption{高重复率不保证有足够的正确匹配。完整设置的十一图统计；接受/正确为每图数量均值，有效图对数只针对正确率。无接受匹配的图对不计入正确率均值，但数量均值包含这些图对的零。定位为重复点中位误差的跨图均值，时间为消融阶段单次计时均值。}
\label{tab:real-full}
\setlength{\tabcolsep}{4pt}
\begin{tabular}{lrrrrrr}
\toprule
方法 & 重复率 & 正确率 & 定位/px & 耗时/ms & 接受/正确 & 有效图对\\
\midrule
Sobel & 0.716 & 0.287 & 2.185 & 362.2 & 1.91/0.73 & 9/11\\
LoG & 0.749 & 0.502 & 2.215 & 355.2 & 2.73/1.64 & 9/11\\
Harris & 0.709 & 0.454 & 2.268 & 391.0 & 2.91/1.55 & 6/11\\
Shi-Tomasi & 0.706 & 0.471 & 2.253 & 380.8 & 3.73/2.09 & 8/11\\
SIFT & 0.547 & 0.541 & 2.237 & 803.7 & 4.91/2.73 & 11/11\\
\bottomrule
\end{tabular}
\end{table}

\clearpage
\subsection{标注几何的近似误差}
表~\ref{tab:real-geometry} 将图对的拟合残差单独列出。\filepath{IMG_0258} 与 \filepath{IMG_0269} 的 RMSE 超过五参考工作像素，\filepath{IMG_0265} 超过四像素；书页弯曲与点位误差混合，不能只把这些残差解释成标注错误，也不能忽略它们。较大的留一最大误差说明部分位置存在外推不稳定，按同一平面阈值判分可能出现假阴性。

\begin{table}[H]
\centering\small
\caption{单一平面不能精确描述全部卷曲书页。点数、单应拟合 RMSE 与留一预测误差，单位均为参考工作像素；这些是标注几何一致性，不是检测器的定位误差。}
\label{tab:real-geometry}
\begin{tabular}{lrrrr}
\toprule
场景图 & 点对数 & 拟合 RMSE & 留一中位 & 留一最大\\
\midrule
IMG\_0258 & 23 & 5.54 & 4.27 & 25.79\\
IMG\_0259 & 23 & 1.06 & 0.81 & 4.55\\
IMG\_0261 & 24 & 2.65 & 2.44 & 6.63\\
IMG\_0262 & 24 & 2.75 & 2.22 & 7.32\\
IMG\_0263 & 24 & 2.78 & 2.35 & 7.27\\
IMG\_0264 & 24 & 2.70 & 2.57 & 7.10\\
IMG\_0265 & 20 & 4.59 & 1.89 & 20.67\\
IMG\_0266 & 24 & 2.89 & 2.56 & 7.00\\
IMG\_0267 & 22 & 3.51 & 3.80 & 7.64\\
IMG\_0268 & 20 & 3.43 & 2.18 & 9.46\\
IMG\_0269 & 22 & 5.47 & 3.92 & 26.39\\
\bottomrule
\end{tabular}
\end{table}

\clearpage
\section{组件关闭消融}
关闭外部金字塔的影响并不统一。Sobel、Harris、SIFT 的重复率变化为 $+0.99,+0.50,+3.12$ 个百分点，LoG 和 Shi-Tomasi 为 $-2.32,-0.21$。SIFT 关闭外部金字塔后的正确率提高 $13.55$ 个百分点，但平均正确匹配数从 $2.73$ 变为 $1.45$；比例上升不等于匹配数量更多。

关闭频域滤波时，SIFT 重复率和正确率分别提高 $8.41$、$38.22$ 个百分点。LoG 正确率则下降 $41.83$ 个百分点，平均正确匹配从 $1.64$ 减为 $0.09$。相反方向的变化说明滤波与描述子互相作用；正确率均值的有效图集也不同，表中的差值不是严格同图子集的配对差。

\begin{table}[H]
\centering\small
\caption{关闭不同组件不必然降低全部指标。各关闭设置均值减去完整设置均值，单位为百分点。P、F、N 分别为关闭外部金字塔、频域滤波、共同 NMS；“--”表示关闭后无接受匹配、正确率未定义。最终预算和阈值口径保持固定。}
\label{tab:real-delta}
\setlength{\tabcolsep}{5pt}
\begin{tabular}{lrrrrrr}
\toprule
 & \multicolumn{3}{c}{重复率变化/pp} & \multicolumn{3}{c}{正确率变化/pp}\\
\cmidrule(lr){2-4}\cmidrule(lr){5-7}
方法 & P & F & N & P & F & N\\
\midrule
Sobel & +0.99 & +2.30 & -25.47 & -13.89 & +7.41 & --\\
LoG & -2.32 & -1.12 & -27.77 & -9.05 & -41.83 & --\\
Harris & +0.50 & -2.31 & -11.90 & -17.72 & -45.42 & -45.42\\
Shi-Tomasi & -0.21 & -5.27 & -18.53 & +11.15 & +11.25 & +41.81\\
SIFT & +3.12 & +8.41 & -19.29 & +13.55 & +38.22 & -4.74\\
\bottomrule
\end{tabular}
\end{table}

\clearpage
关闭共同 NMS 后，五种方法重复率均下降。Sobel、LoG 在十一图中都没有通过双向比例检验的接受匹配，不能把其正确率当作零。Shi-Tomasi 的正确率为 $0.889$，但只在三图有接受匹配，每图平均正确数仅 $0.36$，低于完整设置的 $2.09$；其定位误差从 $2.253$ 降到 $1.464$ 像素也只是条件统计改善。重复文字上的近邻候选可能造成描述子相似、比例检验失败；这是与结果相容的解释，而非已独立证明的唯一原因。

\begin{figure}[H]
\centering
\includegraphics[width=\textwidth]{ablation_metrics.pdf}
\caption{消融需要联合比例、误差和运行时间。四种设置的十一图统计，误差棒为图对之间的样本标准差；正确率面板中未定义的点不绘制，不代表零。接受/正确数量见 CSV，不能只看该图中的比例。}
\end{figure}

\clearpage
\section{多尺度与噪声敏感性}
\subsection{外部金字塔深度}
三层设置并非普遍优于单层。下图的横轴是外部深度而非物理拍摄尺度；每个面板保持一种频域设置，五种管线以颜色、标记及线型共同区分。粗层去噪、细节损失与候选预算竞争同时存在，增加深度的变化应按实际曲线解读。
\begin{figure}[H]
\centering
\includegraphics[width=\textwidth,height=0.72\textheight,keepaspectratio]{repeatability_pyramid.pdf}
\caption{外部层数的增加对各管线并非统一收益。深度一至三层、每层分别预处理，最终总预算固定二百二十；LP、HP 表示低通、高通，误差棒为十一图对样本标准差。SIFT 内部尺度空间始终开启。}
\end{figure}

\clearpage
\subsection{受控数字尺度}
下图只改变场景工作图的重采样倍率，参考不变，几何和 ROI 同步换算，误差仍按参考工作像素计。它补充了原生远近拍摄变化，但不等价于改变焦距、真实光学成像或获得新的独立照片。

不滤波时，LoG 在四个倍率下的重复率依次为 $0.538,0.640,0.738,0.750$，缩小后细节损失较明显；Sobel 为 $0.711,0.680,0.739,0.796$，并非严格单调。SIFT 为 $0.548,0.679,0.631,0.601$，在 $0.75$ 倍处最高，不能将数字放大解释为普遍提高尺度稳定性。上述变化同时包含重采样、候选点数量、固定预算与几何判分的影响。
\begin{figure}[H]
\centering
\includegraphics[width=\textwidth,height=0.72\textheight,keepaspectratio]{repeatability_scale.pdf}
\caption{数字重采样检验的是现有 JPEG 管线的尺度敏感性。场景倍率为 $0.5,0.75,1,1.25$，三层、NMS 开启，五预处理分别成面板；横轴不是金字塔层数，误差棒为十一图对样本标准差。}
\end{figure}

\clearpage
\subsection{受控附加噪声}
在高斯低通下，Sobel 的重复率从附加噪声标准差零时的 $0.716$ 变为二十时的 $0.681$，SIFT 从 $0.547$ 变为 $0.481$。Harris 则约为 $0.709$ 到 $0.710$，LoG 中间噪声强度还出现轻微上升。不能据此断言后两种方法对噪声完全不敏感：密集响应点的几何接近、固定预算和近似 $H$ 都影响这个比例，重复率也不衡量描述子是否正确匹配。
\begin{figure}[H]
\centering
\includegraphics[width=\textwidth,height=0.72\textheight,keepaspectratio]{repeatability_noise.pdf}
\caption{噪声曲线并非全部单调，单一重复率不足以证明可靠性。仅场景额外加入高斯噪声，标准差为 $0,5,10,20$ 灰度级，已有拍摄扰动保留；三层、NMS 开启，误差棒为十一图对样本标准差。}
\end{figure}
